\section{Introduction}
\label{sec:introduction}

Software vulnerabilities remain one of the most persistent threats to modern
systems. In 2023, the National Vulnerability Database (NVD) recorded over
28,000 new CVEs~\cite{nvd2023}. Developers frequently copy code from community
platforms such as Stack Overflow — yet research shows that insecure SO snippets
carry \emph{higher} average scores (14 vs.\ 5) and view counts (36,508
vs.\ 18,713) than their secure counterparts~\cite{chen2019reliable}, making
community reputation an unreliable security signal and motivating automated
detection tools.

Automated vulnerability detection has been approached from two directions.
\emph{Text-based} methods treat source code as a sequence of tokens and apply
NLP models; they scale well but lack program semantics~\cite{vulcnn}.
\emph{Graph-based} methods construct Abstract Syntax Trees (ASTs), Control Flow
Graphs (CFGs), or Program Dependency Graphs (PDGs) and apply Graph Neural
Networks (GNNs); they are semantically rich but slower~\cite{regvd,ignnvd}.

An intermediate line of work — starting with VulCNN~\cite{vulcnn} and refined
by VulGAI~\cite{vulgai} — converts the PDG into an image via node-centrality
encoding and uses a CNN for classification. This achieves the speed of
image-recognition models while preserving some structural information, but the
CNN's local receptive field cannot perform graph message-passing and loses the
relational structure embedded in the PDG. VulSCC~\cite{vulscc} augments this
with a code LLM, but still lacks the graph branch. In the smart-contract domain,
TMF-Net~\cite{tmfnet} fuses graph, image, and bytecode modalities — but targets
Solidity contracts, not general C/C++/Java source code.

Prior work such as TemCon~\cite{temcon} demonstrates that structural code 
patterns are informative for identifying and fixing concurrency defects. VulGCL
extends this insight to security vulnerabilities, combining structural, visual, and
semantic representations for detection. 

\textbf{Gap.} No prior work simultaneously fuses (1) graph structural features,
(2) visual code representations, and (3) LLM semantic embeddings within a
single end-to-end vulnerability detector for general source code.

We address this gap with \textbf{VulGCL}, which makes the following
contributions:

\begin{itemize}
  \item \textbf{Multimodal architecture.} VulGCL is the first framework to
    jointly learn graph (GNN), visual (CNN), and semantic (CodeBERT) features
    for general source-code vulnerability detection~(\S\ref{sec:methodology}).

  \item \textbf{Community-signal labeling.} We construct a supplementary
    training corpus from Stack Overflow code snippets labeled via static
    analysis, using comment-flag patterns and accepted-answer status as weak
    supervision signals — informed by the finding that raw vote counts
    anticorrelate with security~\cite{chen2019reliable}~(\S\ref{sec:data}).

  \item \textbf{Edge deployment.} VulGCL-Lite, a quantized variant of VulGCL,
    runs real-time inference on a Raspberry Pi 4B, enabling vulnerability
    scanning on IoT and embedded systems~(\S\ref{sec:deployment}).

  \item \textbf{Empirical evaluation.} Extensive experiments on Devign and
    BigVul show VulGCL outperforms all single-modality and two-modality
    baselines~(\S\ref{sec:experiments}).
\end{itemize}
